% Created 2024-04-04 jeu. 10:26
% Intended LaTeX compiler: pdflatex
\documentclass[11pt]{article}
\usepackage[utf8]{inputenc}
\usepackage[T1]{fontenc}
\usepackage{graphicx}
\usepackage{longtable}
\usepackage{wrapfig}
\usepackage{rotating}
\usepackage[normalem]{ulem}
\usepackage{amsmath}
\usepackage{amssymb}
\usepackage{capt-of}
\usepackage{hyperref}
\date{\today}
\title{Plan détaillé du manuscript}
\hypersetup{
 pdfauthor={},
 pdftitle={Plan détaillé du manuscript},
 pdfkeywords={},
 pdfsubject={},
 pdfcreator={Emacs 29.3 (Org mode 9.7)}, 
 pdflang={French}}
\usepackage{biblatex}

\begin{document}

\maketitle
\tableofcontents

\subsection{Intro}
\label{sec:orgea75187}
\section{Conception d'un système de réévaluation basé sur les règles de transformation de graphe}
\label{sec:org66fccae}
\subsection{Un état de l'art}
\label{sec:org227c66d}
\subsubsection{rejeu}
\label{sec:org542bf54}
\begin{enumerate}
\item modélisation paramétrique
\label{sec:orgf995bd1}
\item réévaluation (nommage persistant entre autre) : kripac, capoyleas, chen, mun, wu, cheon, fajarna, agbodan, marcheix, hepworth,…
\label{sec:orgf5d77c2}
\end{enumerate}
\subsubsection{modélisation à base de règles}
\label{sec:orgd6b046f}
\begin{enumerate}
\item G-carte
\label{sec:org105037a}
\begin{itemize}
\item G-carte
\item Orbite
\end{itemize}
\item Jerboa
\label{sec:org528b1d0}
\begin{itemize}
\item Règles de transformation de graphe
\item Règles Jerboa de transformation de graphe
\item Chemin (?)
\item Lexique pour les G-cartes
\begin{itemize}
\item brin, liaisons, orbites
\end{itemize}
\item Lexique pour les règles
\begin{itemize}
\item nœud, arcs, orbites
\end{itemize}
\end{itemize}
\item Rejeu à base de règle (Anaïs)
\label{sec:orge1dc033}
\end{enumerate}
\subsection{Formalisation des évènements liés à l'application des opérations de modélisation}
\label{sec:org629a912}
\subsubsection{Un état des lieux de la détection d'évènement}
\label{sec:org5ad7613}
\subsubsection{Les évènements dans un graphe/une G-carte}
\label{sec:org64acdc5}
\begin{enumerate}
\item Formaliser la création
\label{sec:org9ea4717}
\item Formaliser la suppression
\label{sec:org9f42cfb}
\item Formaliser le non-changement
\label{sec:org65ba127}
\item Formaliser la modification
\label{sec:orgee3c261}
\item Formaliser la scission
\label{sec:org1f03155}
\item Formaliser la fusion
\label{sec:org49689fe}
\item (Formaliser le sans effet ?)
\label{sec:org3102c7b}
\end{enumerate}
\subsubsection{Les évènements dans les règles de transformation de graphe}
\label{sec:org8382405}
\begin{enumerate}
\item Formaliser la création
\label{sec:orgd3d3d2a}
\item Formaliser la suppression
\label{sec:org017dd76}
\item Formaliser le non-changement
\label{sec:org83b4051}
\item Formaliser la modification
\label{sec:org21c2105}
\item Formaliser la scission
\label{sec:org879d2a1}
\item Formaliser la fusion
\label{sec:orgbdc1462}
\item (Formaliser le sans effet ?)
\label{sec:org0cf3004}
\end{enumerate}
\subsubsection{Les évènements dans les règles Jerboa}
\label{sec:orgdff4c4e}
\begin{enumerate}
\item Formaliser la création
\label{sec:orge50a91d}
\item Formaliser la suppression
\label{sec:org8165c48}
\item Formaliser le non-changement
\label{sec:orgd8a60e9}
\item Formaliser la modification
\label{sec:orgdd49ef2}
\item Formaliser la scission
\label{sec:org6cc50a3}
\item Formaliser la fusion
\label{sec:orgbbdb643}
\item (Formaliser le sans effet ?)
\label{sec:org7c756cb}
\end{enumerate}
\subsubsection{Application : suivi de l'évolution et de l'origine d'une orbite}
\label{sec:org95820e9}
\begin{enumerate}
\item Formaliser l'origine d'un évènement créé par une règle de transformation de graphe
\label{sec:orgc3a6b82}
\item Formaliser l'origine d'un évènement créé par une règle Jerboa
\label{sec:orge387459}
\end{enumerate}
\subsection{Réévaluation dans les systèmes de modélisation à base de règles de transformation de graphe}
\label{sec:org7d14af0}
\subsubsection{Conception d'un système de spécification paramétrique}
\label{sec:org71f29e0}
\begin{enumerate}
\item La spécification paramétrique : un enregistrement du processus de construction d'un modèle
\label{sec:org39b7fd8}
\begin{enumerate}
\item Une identification des opérations
\label{sec:orgfd8d64c}
\item Des paramètres géométriques
\label{sec:orgf7ec21a}
\item Des paramètres topologiques
\label{sec:org74ffa3d}
\end{enumerate}
\item Le nommage persistant : un mécanisme pour identifier et suivre les entités topologiques
\label{sec:orgf32fb09}
\begin{enumerate}
\item Mise en place d'un mécanisme de nommage à base de règles de transformation de graphe
\label{sec:org5fa9adc}
\end{enumerate}
\item Un système de nommage persistant basé sur les règles Jerboa
\label{sec:orgfe2f539}
Les règles, mises bout-à-bout, nous permettent de tracer l'historique des
évènements impliqué dans la création d'un brin
\end{enumerate}
\subsubsection{DAGs évaluation}
\label{sec:org3ea6401}
\begin{enumerate}
\item Construction à partir d'un nom persistant (de la fin vers le début)
\label{sec:org165efbe}
\item Les origines pour un historique complet des cellules topologiques
\label{sec:org95918f8}
Combinaison d'orbites d'origine et de suivi dans un DAG pour une représentation
complète d'un évènement entre deux applications de règles
\item Des chemins pour plus de précision
\label{sec:org190303a}
Formalisation de l'usage de chemins pour préciser une origine lorsque celle-ci
n'est pas directement liée à un nœud d'accroche dans une règle
\end{enumerate}
\subsubsection{DAGs réévaluation}
\label{sec:orga954360}
\begin{enumerate}
\item Construction à partir d'un DAG d'évaluation (du début jusqu'à la fin)
\label{sec:org02d5b93}
\item La réévaluation, un contexte différent
\label{sec:org039e316}
Puisque les DAGs de réévaluation sont utilisés pour l'appariement de paramètres
topologiques, ceux-ci ne contiennent non pas des nœuds de règles mais des brins
issus de l'objet en cours de reconstruction.
\item Un mécanisme extensible via des stratégies de réévaluation
\label{sec:orgd4e36d0}
Des stratégies par défaut
\begin{itemize}
\item Scission d'un paramètre topologique : évaluation de tous les paramètres éligibles
\item Suppression d'un paramètre topologique : interruption de la réévaluation pour ce paramètres (potentiellement de l'objet)
\end{itemize}
\end{enumerate}
\subsection{Réévaluer des scripts de règles pour générer des modèles complexes (En cours)}
\label{sec:orga3ae81b}
\subsubsection{Les scripts Jerboa}
\label{sec:org9689965}
\subsubsection{support des scripts}
\label{sec:orgcc5b6ea}
\begin{enumerate}
\item Les séquences de règles
\label{sec:org34ee3bb}
\begin{itemize}
\item Extension de la spécification paramétrique : des blocs « séquence » de règles contenues dans un script
\item Extension des DAGs : une boîte englobante contenant les règles jouées dans le contexte d'un script
\end{itemize}
\item Les boucles de sous-orbites
\label{sec:org19594ec}
\begin{itemize}
\item Extension de la spécification paramétrique : des blocs « boucle » de règles contenues dans une boucle de script
\item Extension des DAGs : Ajout de branches pour identifier les paramètres topologiques itérés dans une boucle
\end{itemize}
\item Les conditions Si-Sinon-Alors
\label{sec:org08bc9aa}
\begin{itemize}
\item Extension de la spécification paramétrique : des blocs « condition » de règles contenues dans une condition de script
\item Extension des DAGs : Ajout d'un paramètre associé à boîte permettant d'identifier une instance de cellule topologique (spécifique aux règles Jerboa)
\end{itemize}
\end{enumerate}
\subsection{Implantation d'un système de modélisation à base de règles de transformation de graphe}
\label{sec:orga1fe187}
\subsubsection{Structure et implantation d'une spécification paramétrique}
\label{sec:org8762ecd}
\subsubsection{Structure et implantation du nommage persistant}
\label{sec:org093133a}
\subsubsection{Structure et implantation des DAGs de jeu et de rejeu}
\label{sec:orgf314be9}
\subsubsection{Implantation d'un outil pour visualiser les DAGs}
\label{sec:orgb519670}
\subsubsection{Implantation d'une Interface Utilisateur pour éditer et générer des spécifications paramétriques (À boucler absolument)}
\label{sec:org33eb70c}
\section{Animation spatio-temporelle (Actuellement, totalement de la perspective)}
\label{sec:orgcb98af8}
\subsection{Un état de l'art}
\label{sec:org5648a0d}
\begin{itemize}
\item de luca, stefani,…
\end{itemize}
\subsection{Extension de la spécification paramétrique pour annoter avec des données temporelles}
\label{sec:org88ad9d0}
\subsubsection{Annotation des opérations avec des données temporelles}
\label{sec:org3134710}
\subsubsection{Modulation de la vitesse d'animation en fonction des deltas temporels entre deux opérations}
\label{sec:org83180aa}
\subsection{Processus d'animation de modèle générés automatiquement}
\label{sec:org6e3ee11}
\subsubsection{Génération à la volée des modèles ?}
\label{sec:org5d8e99e}
\subsubsection{Pré-génération des modèles avant phase d'animation ?}
\label{sec:org5def7fa}
\end{document}
